\documentclass[11pt]{article}

\usepackage[margin=1in]{geometry}
\usepackage{amsmath,amssymb,amsthm}
\usepackage{array}
\emergencystretch=1.5em

\newtheorem{theorem}{Theorem}[section]
\newtheorem{proposition}[theorem]{Proposition}
\newtheorem{lemma}[theorem]{Lemma}
\newtheorem{corollary}[theorem]{Corollary}
\theoremstyle{definition}
\newtheorem{definition}[theorem]{Definition}
\newtheorem{example}[theorem]{Example}
\theoremstyle{remark}
\newtheorem{remark}[theorem]{Remark}

\newcommand{\rk}{\mathrm{rk}}
\newcommand{\hatzero}{\hat{0}}
\newcommand{\hatone}{\hat{1}}
\newcommand{\Z}{\mathbb{Z}}
\newcommand{\N}{\mathbb{N}}
\newcommand{\Q}{\mathbb{Q}}

\title{\textbf{Resolution of Conjecture 00000008539:\\
M\"obius Function Absolute Values, Log-Concavity,\\
and Global Equality on Geometric Lattices}}
\author{Formal Disproof, Complete Classification, and Core Lean 4 Verification}
\date{}

\begin{document}
\maketitle

\begin{abstract}
We resolve Conjecture \texttt{00000008539} from \emph{The Last Math Competition}, which asserts that for a geometric lattice $L$ with incidence-algebra M\"obius function $\mu = \zeta^{-1}$, the absolute values of the M\"obius function form a log-concave sequence with first term $1$, with global equality holding only for Boolean lattices. We analyze every natural mathematical convention by which the absolute values of $\mu$ yield a rank-indexed sequence on a finite geometric lattice $L$ of rank $n$, and establish that the conjecture as stated is \textbf{false} under each of them, while proving sharp structural classifications explaining both halves of the statement:
\begin{enumerate}
  \item \textbf{Unsigned Whitney numbers of the first kind} $w_k^+(L) = \sum_{\rk(x)=k} |\mu(\hatzero, x)|$: While $w_0^+(L)=1$ and log-concavity $(w_k^+)^2 \ge w_{k-1}^+ w_{k+1}^+$ hold for every finite geometric lattice (Adiprasito--Huh--Katz, 2018), global equality \emph{fails} on every Boolean lattice $B_n$ of rank $n \ge 2$ (already $w^+(B_2)=(1,2,1)$ gives $2^2 = 4 > 1 \cdot 1$). Stronger still, we prove from Rota's sign theorem and $\sum_{x \in L} \mu(\hatzero,x)=0$ that \emph{no} finite geometric lattice of rank $n \ge 2$ ever satisfies global equality; thus global equality holds if and only if $\rk(L) \le 1$ (vacuously over an empty interior index set).
  \item \textbf{Elementwise, chain, rank-averaged, and rank-extremal $|\mu|$ sequences}: Along any maximal chain $\hatzero = x_0 \lessdot \dots \lessdot x_n = \hatone$ (or rank-selector, rank-average $w_k^+(L)/W_k(L)$, rank-minimum, or rank-maximum), every Boolean lattice $B_n$ yields $(1,1,\dots,1)$, attaining global equality for all $n \ge 0$. However, we prove $|\mu(\hatzero,\hatone)| = 1$ if and only if $L$ is Boolean, and deduce that log-concavity \emph{fails} on \emph{every} non-Boolean geometric lattice (minimal counterexample $M_3$: $(1,1,2)$, where $1^2 < 1 \cdot 2$).
  \item \textbf{Binomial-normalized Whitney numbers} $b_k(L) = w_k^+(L)/\binom{n}{k}$ (ultra-log-concavity): Even after normalizing so that $B_n$ gives $(1,\dots,1)$, log-concavity fails on the rank-$3$ uniform geometric lattice $U_{3,4}$, where $b(U_{3,4})=(1, 4/3, 2, 3)$ violates $(4/3)^2 \ge 1 \cdot 2$.
\end{enumerate}
All definitions, counterexamples, non-vacuity witnesses, Weisner's identity, Rota's sign theorem, the Boolean characterization $|\mu(\hatzero,\hatone)|=1$, and the general classification theorems are mechanized from first principles in self-contained Core Lean 4 (\texttt{v4.19.0}, zero external dependencies, zero \texttt{sorry}) and independently verified in Python across $14$ geometric lattices.
\end{abstract}

\section{Problem Statement and Overview of Conventions}

\subsection{Original Statement}
Conjecture \texttt{00000008539} (\texttt{SOURCE.md}) states:
\begin{quote}
\textbf{Definition:} The Mobius function is the inverse of zeta in the incidence algebra.\\
\textbf{Conjecture:} The absolute values of the Mobius function of a geometric lattice form a log-concave sequence with first term 1; and equality holds globally only for Boolean lattices. (\emph{Mobius log-concavity})
\end{quote}
The accompanying Chinese formulation (\emph{``jihege de M\"obius jueduizhi goucheng duishu ao xulie qie shouxiang wei 1; denghao jin zai buerge zhengti chengli''}) carries the identical mathematical structure, using the standard equality-case phrasing \emph{``equality holds globally [across all interior indices simultaneously] only at Boolean lattices''}.

\subsection{Why Multi-Convention Analysis Is Essential}
For a finite geometric lattice $L$ of rank $n \ge 0$ with minimum $\hatzero$, maximum $\hatone$, and rank function $\rk : L \to \{0,1,\dots,n\}$, the M\"obius function $\mu(\hatzero, x)$ is indexed by elements $x \in L$ rather than directly by integers $k \in \{0,1,\dots,n\}$. Turning ``the absolute values of the M\"obius function'' into a sequence $(a_0, a_1, \dots, a_n)$ admits three natural mathematical conventions:
\begin{itemize}
  \item \textbf{Convention I (Rank-summed / Unsigned Whitney numbers of the first kind):}
  \[
    w_k^+(L) \;=\; \sum_{\substack{x \in L \\ \rk(x)=k}} |\mu(\hatzero, x)| \qquad (0 \le k \le n).
  \]
  In matroid and poset theory, ``M\"obius log-concavity'' refers to the Heron--Rota--Welsh conjecture for $w_k^+(L)$, proved by Adiprasito, Huh, and Katz~\cite{AHK2018}.
  \item \textbf{Convention II (Elementwise along maximal chains / rank-selectors, or rank-averaged / extremal):}
  \[
    a_k \;=\; |\mu(\hatzero, x_k)| \quad (\rk(x_k)=k), \qquad \text{or} \qquad \bar{\mu}_k(L) \;=\; \frac{w_k^+(L)}{W_k(L)},
  \]
  where $W_k(L) = |\{x \in L : \rk(x)=k\}|$ is the $k$-th Whitney number of the second kind.
  \item \textbf{Convention III (Binomial-normalized Whitney numbers / Ultra-log-concavity):}
  \[
    b_k(L) \;=\; \frac{w_k^+(L)}{\binom{n}{k}} \qquad (0 \le k \le n).
  \]
\end{itemize}

Table~\ref{tab:summary} summarizes the fate of every clause of Conjecture \texttt{00000008539} across all conventions. Notice the structural diagnosis of the conjecture: the author combined the log-concavity of \emph{rank-summed} Whitney numbers $w_k^+(L)$ (Convention I) with the global constant-$(1,\dots,1)$ equality behavior on Boolean lattices $B_n$ belonging to \emph{elementwise / normalized} $|\mu|$ (Conventions II--III). Under any single consistent definition, the conjecture is \textbf{false}.

\begin{table}[ht]
\centering
\small
\setlength{\tabcolsep}{4.5pt}
\renewcommand{\arraystretch}{1.28}
\begin{tabular}{@{}>{\raggedright\arraybackslash}p{3.95cm}ccc>{\raggedright\arraybackslash}p{3.65cm}@{}}
\hline
\textbf{Sequence Convention} & $a_0 = 1$ & \textbf{Log-Concave} & \textbf{Eq.\ $\Leftrightarrow$ Bool} & \textbf{Counterexample} \\
\hline
\textbf{I.} Unsigned Whitney $w_k^+(L)$ & True & True~\cite{AHK2018} & \textbf{False} ($n \ge 2$) & $B_2$: $(1,2,1)$, $2^2 \neq 1$ \\
\textbf{IIA.} Chain $|\mu(\hatzero, x_k)|$ & True & \textbf{False} (non-Bool) & True on $B_n$ & $M_3$: $(1,1,2)$, $1^2 < 2$ \\
\textbf{IIB.} Average $w_k^+(L)/W_k(L)$ & True & \textbf{False} (non-Bool) & True on $B_n$ & $M_3$: $(1,1,2)$, $1^2 < 2$ \\
\textbf{IIC.} Extremal $\min/\max |\mu|$ & True & \textbf{False} (non-Bool) & True on $B_n$ & $M_3$: $(1,1,2)$, $1^2 < 2$ \\
\textbf{III.} Normalized $w_k^+(L)/\binom{n}{k}$ & True & \textbf{False} & True on $B_n$ & $U_{3,4}$: $(1,\tfrac{4}{3},2,3)$, $\tfrac{16}{9} < 2$ \\
\hline
\end{tabular}
\caption{Truth values and minimal counterexamples across all natural conventions.}
\label{tab:summary}
\end{table}

\section{Incidence Algebras and Geometric Lattices}

We record the exact definitions formalized in \texttt{lean/Main.lean}.

\begin{definition}[Incidence Algebra and M\"obius Function]
Let $(P, \le)$ be a finite partially ordered set. The \emph{incidence algebra} $I(P, \Z)$ consists of all functions $f : P \times P \to \Z$ such that $f(x, y) = 0$ whenever $x \not\le y$, equipped with pointwise addition and convolution
\[
  (f * g)(x, y) \;=\; \sum_{z \in P} f(x, z)\,g(z, y) \;=\; \sum_{x \le z \le y} f(x, z)\,g(z, y).
\]
The multiplicative identity is $\delta(x, y) = [x = y]$, and the \emph{zeta function} is $\zeta(x, y) = [x \le y]$. An element $\mu \in I(P, \Z)$ is the \emph{M\"obius function} of $P$ if $\mu * \zeta = \delta$ and $\zeta * \mu = \delta$.
\end{definition}

\begin{proposition}[Existence, Two-Sided Invertibility, and Uniqueness of $\mu$]
\label{prop:mobius}
For every finite poset $(P, \le)$, $I(P, \Z)$ is an associative unital ring under convolution, and $\zeta$ admits a unique two-sided inverse $\mu \in I(P, \Z)$, recursively determined for $x \le y$ by
\[
  \mu(x, x) = 1, \qquad \mu(x, y) = -\sum_{x \le z < y} \mu(x, z) \quad (x < y).
\]
In particular, if $y$ covers $x$ ($x \lessdot y$, meaning $x < y$ and $\{z \in P : x \le z \le y\} = \{x, y\}$), then $\mu(x, y) = -1$.
\end{proposition}

\begin{proof}
Associativity of $*$ follows by Fubini summation over $P$, and $\delta * f = f = f * \delta$ holds because $P$ is finite. Defining $\mu_{\mathrm{L}}(x,y)$ by strong recursion on $|[x,y]|$ via $\mu_{\mathrm{L}}(x,x)=1$ and $\mu_{\mathrm{L}}(x,y) = -\sum_{x \le z < y} \mu_{\mathrm{L}}(x,z)$ for $x<y$ immediately yields $\mu_{\mathrm{L}} * \zeta = \delta$. Symmetrically, defining $\mu_{\mathrm{R}}(x,y)$ by $\mu_{\mathrm{R}}(x,x)=1$ and $\mu_{\mathrm{R}}(x,y) = -\sum_{x < z \le y} \mu_{\mathrm{R}}(z,y)$ yields $\zeta * \mu_{\mathrm{R}} = \delta$. By associativity,
\[
  \mu_{\mathrm{L}} \;=\; \mu_{\mathrm{L}} * \delta \;=\; \mu_{\mathrm{L}} * (\zeta * \mu_{\mathrm{R}}) \;=\; (\mu_{\mathrm{L}} * \zeta) * \mu_{\mathrm{R}} \;=\; \delta * \mu_{\mathrm{R}} \;=\; \mu_{\mathrm{R}}.
\]
Thus $\mu = \mu_{\mathrm{L}}$ satisfies both $\mu * \zeta = \delta$ and $\zeta * \mu = \delta$, and any $m$ with $m * \zeta = \delta$ satisfies $m = m * (\zeta * \mu) = (m * \zeta) * \mu = \mu$. When $x \lessdot y$, the interval $[x,y]$ is $\{x,y\}$, so $\mu(x,y) = -\mu(x,x) = -1$.
\end{proof}

\begin{definition}[Geometric Lattices and Boolean Lattices]
Let $(L, \le, \hatzero, \hatone, \wedge, \vee)$ be a finite bounded lattice. An element $p \in L$ is an \emph{atom} if $\hatzero \lessdot p$.
\begin{enumerate}
  \item $L$ is a \emph{finite geometric lattice} (Birkhoff; Stanley~\cite[\S 3.3]{StanleyEC1}) if:
  \begin{enumerate}
    \item (\emph{Atomistic}) Every element is the join of the atoms below it: for all $x, y \in L$, if every atom $p \le x$ satisfies $p \le y$, then $x \le y$.
    \item (\emph{Graded and submodular}) There is a rank function $\rk : L \to \N$ satisfying $\rk(\hatzero) = 0$, $\rk(y) = \rk(x) + 1$ whenever $x \lessdot y$, and $\rk(x \vee y) + \rk(x \wedge y) \le \rk(x) + \rk(y)$ for all $x, y \in L$.
  \end{enumerate}
  We write $\mathrm{rank}(L) = \rk(\hatone)$. Note that $\rk$ is uniquely determined by $(L, \le)$: in any finite poset, every relation $\hatzero < y$ admits a saturated chain $\hatzero = y_0 \lessdot y_1 \lessdot \dots \lessdot y_k = y$, forcing $\rk(y) = k$. Moreover, $p \in L$ is an atom if and only if $\rk(p) = 1$, and every rank $k \in \{0,\dots,\mathrm{rank}(L)\}$ is inhabited.
  \item $L$ is a \emph{Boolean lattice} if it is \emph{distributive} ($x \wedge (y \vee z) = (x \wedge y) \vee (x \wedge z)$ for all $x,y,z \in L$) and \emph{complemented} (for every $x \in L$, there exists $y \in L$ with $x \wedge y = \hatzero$ and $x \vee y = \hatone$).
\end{enumerate}
\end{definition}

\begin{definition}[Log-Concavity and Global Equality]
A sequence $(a_0, a_1, \dots, a_n)$ of positive real (or rational/integer) numbers has \emph{first term $1$} if $a_0 = 1$, is \emph{log-concave} if $a_k^2 \ge a_{k-1} a_{k+1}$ for all $1 \le k \le n - 1$ (since $a_k > 0$, internal zeros never arise), and satisfies \emph{global equality} if $a_k^2 = a_{k-1} a_{k+1}$ for all $1 \le k \le n - 1$.
\end{definition}

\section{Weisner's Theorem, Rota's Sign Theorem, and Boolean Characterization}

To prove universal structural theorems rather than merely isolated counterexamples, we establish Weisner's identity, Rota's strict sign theorem, and the M\"obius characterization of Boolean geometric lattices directly from the axioms above.

\begin{theorem}[Weisner's Identity]
\label{thm:weisner}
Let $L$ be any finite lattice, let $a \in L$ satisfy $\hatzero < a$, and let $y \in L$. Then
\[
  \sum_{\substack{x \in L \\ x \vee a = y}} \mu(\hatzero, x) \;=\; 0.
\]
\end{theorem}

\begin{proof}
For every $z \in L$ with $a \le z \le y$, we have $\hatzero < a \le z$, hence $\hatzero \neq z$ and $(\mu * \zeta)(\hatzero, z) = \delta(\hatzero, z) = 0$. Weighting by $\mu(z, y)$ and summing over $a \le z \le y$ gives
\[
  0 \;=\; \sum_{a \le z \le y} \left(\sum_{\hatzero \le x \le z} \mu(\hatzero, x)\right) \mu(z, y).
\]
For $x, z \in L$, the pair of inequalities $(x \le z \text{ and } a \le z \le y)$ is equivalent by the universal property of join to $(x \le y \text{ and } x \vee a \le z \le y)$. Exchanging the finite sums and applying $\zeta * \mu = \delta$ (Proposition~\ref{prop:mobius}) yields
\[
  0 \;=\; \sum_{\hatzero \le x \le y} \mu(\hatzero, x) \left(\sum_{x \vee a \le z \le y} \mu(z, y)\right) \;=\; \sum_{\hatzero \le x \le y} \mu(\hatzero, x)\,\delta(x \vee a, y) \;=\; \sum_{\substack{x \in L \\ x \vee a = y}} \mu(\hatzero, x),
\]
since $x \vee a = y$ automatically implies $\hatzero \le x \le y$.
\end{proof}

\begin{lemma}[Existence of Relative Coatom Complements]
\label{lem:coatom}
Let $L$ be a finite geometric lattice, let $p \in L$ be an atom, and let $y, z \in L$ satisfy $p \le y$, $z \le y$, and $p \not\le z$. Then there exists $x \in L$ such that $z \le x$, $p \not\le x$, $x \vee p = y$, and $x \neq y$. Furthermore, every $x \neq y$ satisfying $x \vee p = y$ has $\rk(y) = \rk(x) + 1$.
\end{lemma}

\begin{proof}
The finite set $S = \{w \in L : z \le w \le y \text{ and } p \not\le w\}$ is non-empty because $z \in S$. Choose $x \in S$ maximizing $\rk(x)$. Since $p \not\le x$ and $p \le y$, we have $x \neq y$ and $x \vee p \le y$. Suppose for contradiction that $x \vee p < y$. Since $L$ is atomistic and $y \not\le x \vee p$, there exists an atom $q \le y$ with $q \not\le x \vee p$; in particular, $q \not\le x$. Because $q$ is an atom ($\hatzero \lessdot q$) and $q \not\le x$, we have $x \wedge q = \hatzero$, so submodularity gives $\rk(x \vee q) \le \rk(x) + \rk(q) - \rk(\hatzero) = \rk(x) + 1$. Since $x < x \vee q$, strict monotonicity of $\rk$ yields $\rk(x \vee q) = \rk(x) + 1$. Identically, $x \wedge p = \hatzero$ and $\rk(x \vee p) = \rk(x) + 1$. Now $z \le x < x \vee q \le y$ and $\rk(x \vee q) > \rk(x)$, so maximality of $\rk(x)$ on $S$ forces $x \vee q \notin S$, hence $p \le x \vee q$. Therefore $x \vee p \le x \vee q$ while $\rk(x \vee p) = \rk(x) + 1 = \rk(x \vee q)$, forcing $x \vee p = x \vee q \ge q$, contradicting $q \not\le x \vee p$. Thus $x \vee p = y$. Finally, for any $x \neq y$ with $x \vee p = y$, we have $p \not\le x$ (otherwise $y = x \vee p = x$), so $x \wedge p = \hatzero$, $x < y$, and submodularity gives $\rk(x) < \rk(y) \le \rk(x) + 1$, hence $\rk(y) = \rk(x) + 1$.
\end{proof}

\begin{theorem}[Rota's Strict Sign Theorem~\cite{Rota1964}]
\label{thm:rota}
For every finite geometric lattice $L$ and every $y \in L$,
\[
  (-1)^{\rk(y)} \mu(\hatzero, y) \;\ge\; 1, \qquad \text{so} \qquad |\mu(\hatzero, y)| \;=\; (-1)^{\rk(y)} \mu(\hatzero, y) \;\ge\; 1.
\]
\end{theorem}

\begin{proof}
We proceed by strong induction on $\rk(y)$. If $\rk(y) = 0$, then $y = \hatzero$ and $(-1)^0 \mu(\hatzero, \hatzero) = 1 \ge 1$. If $\rk(y) = k + 1 \ge 1$, then $y \neq \hatzero$, so by atomisticity there exists an atom $p \le y$. Separating the $x = y$ term in Theorem~\ref{thm:weisner} (with $a = p$) and multiplying by $(-1)^{k+1}$, Lemma~\ref{lem:coatom} ($\rk(x) = k$) gives
\begin{equation}
\label{eq:signed-weisner}
  (-1)^{\rk(y)} \mu(\hatzero, y) \;=\; \sum_{\substack{x \in L \\ x \vee p = y,\; x \neq y}} (-1)^{\rk(x)} \mu(\hatzero, x).
\end{equation}
By the induction hypothesis, each summand on the right-hand side is $\ge 1$, and applying Lemma~\ref{lem:coatom} with $z = \hatzero$ guarantees at least one summand exists. Hence $(-1)^{\rk(y)} \mu(\hatzero, y) \ge 1$.
\end{proof}

\begin{theorem}[M\"obius Characterization of Boolean Geometric Lattices]
\label{thm:bool-char}
Let $L$ be a finite geometric lattice. Then $L$ is a Boolean lattice if and only if $|\mu(\hatzero, \hatone)| = 1$. Moreover, when $L$ is Boolean, $|\mu(\hatzero, y)| = 1$ for every $y \in L$.
\end{theorem}

\begin{proof}
($\Rightarrow$) Suppose $L$ is Boolean (in particular distributive). We prove $(-1)^{\rk(y)} \mu(\hatzero, y) = 1$ for all $y \in L$ by strong induction on $\rk(y)$. The base case $y = \hatzero$ is trivial. For $y > \hatzero$, pick an atom $p \le y$ and an element $x_0 \neq y$ with $x_0 \vee p = y$ via Lemma~\ref{lem:coatom}. For any $x \neq y$ with $x \vee p = y$, we have $p \not\le x$ and hence $x \wedge p = \hatzero = x_0 \wedge p$. By distributivity,
\[
  x \;=\; x \wedge y \;=\; x \wedge (x_0 \vee p) \;=\; (x \wedge x_0) \vee (x \wedge p) \;=\; x \wedge x_0 \;=\; x_0 \wedge x \;=\; x_0 \wedge (x \vee p) \;=\; x_0.
\]
Thus~\eqref{eq:signed-weisner} has the single summand $x = x_0$, giving $(-1)^{\rk(y)} \mu(\hatzero, y) = (-1)^{\rk(x_0)} \mu(\hatzero, x_0) = 1$.

($\Leftarrow$) Conversely, suppose $|\mu(\hatzero, \hatone)| = 1$, so $(-1)^{\rk(\hatone)} \mu(\hatzero, \hatone) = 1$ by Theorem~\ref{thm:rota}. For every atom $p \in L$, applying~\eqref{eq:signed-weisner} at $y = \hatone$ writes $1$ as a sum of integers each $\ge 1$ (Theorem~\ref{thm:rota}) over a non-empty set (Lemma~\ref{lem:coatom} with $z = \hatzero$). Consequently, there is a \emph{unique} element $c_p \in L$ satisfying $c_p \vee p = \hatone$ and $c_p \neq \hatone$. Note that $p \not\le c_p$. Now for any $z \in L$: if $z \le c_p$, then $p \not\le z$; conversely, if $p \not\le z$, Lemma~\ref{lem:coatom} (applied to $y = \hatone$ and $z$) produces $x \ge z$ with $x \vee p = \hatone$ and $x \neq \hatone$, so uniqueness forces $x = c_p$ and hence $z \le c_p$. Thus
\[
  \forall \text{ atoms } p \in L,\; \forall z \in L, \qquad p \not\le z \;\Longleftrightarrow\; z \le c_p.
\]
Distributivity follows immediately: in any lattice $(x \wedge y) \vee (x \wedge z) \le x \wedge (y \vee z)$; conversely, if an atom $p$ satisfies $p \le x$ and $p \le y \vee z$, we cannot have both $p \not\le y$ and $p \not\le z$ (else $y \le c_p$ and $z \le c_p$, giving $y \vee z \le c_p$, so $p \not\le y \vee z$), hence $p \le x \wedge y$ or $p \le x \wedge z$. By atomisticity, $x \wedge (y \vee z) \le (x \wedge y) \vee (x \wedge z)$. Finally, for any $x \in L$, let $y = \bigvee \{p \text{ atom} : p \not\le x\}$. Every atom lies below $x$ or $y$, so $x \vee y = \hatone$ by atomisticity; if an atom $q$ satisfied $q \le x \wedge y$, then every atom $p$ with $p \not\le x$ satisfies $q \neq p$, hence $q \not\le p$, hence $p \le c_q$, forcing $y \le c_q$ and $q \not\le y$, a contradiction. Thus $x \wedge y = \hatzero$, so $L$ is complemented and distributive.
\end{proof}

\begin{lemma}[Geometric Lattices of Rank $\le 1$]
\label{lem:rk1-bool}
Every finite geometric lattice $L$ with $\mathrm{rank}(L) \le 1$ is a Boolean lattice ($L \cong B_0$ or $L \cong B_1$).
\end{lemma}

\begin{proof}
For every $x \in L$, $0 \le \rk(x) \le \rk(\hatone) \le 1$. If $x \neq \hatzero$, strict monotonicity of $\rk$ gives $1 \le \rk(x)$, and if also $x \neq \hatone$, then $1 \le \rk(x) < \rk(\hatone) \le 1$, impossible. Hence $L = \{\hatzero, \hatone\}$, for which distributivity and complementation ($\hatzero \leftrightarrow \hatone$) hold by direct case check.
\end{proof}

\section{Convention I: Unsigned Whitney Numbers of the First Kind}

Consider first the standard sequence of unsigned Whitney numbers of the first kind,
\[
  w_k^+(L) \;=\; \sum_{\substack{x \in L \\ \rk(x)=k}} |\mu(\hatzero, x)| \qquad (0 \le k \le n = \mathrm{rank}(L)).
\]
Clearly $w_0^+(L) = |\mu(\hatzero, \hatzero)| = 1$, and by the theorem of Adiprasito, Huh, and Katz~\cite{AHK2018}, the sequence $w^+(L) = (w_0^+(L), \dots, w_n^+(L))$ is log-concave for every finite geometric lattice $L$. Nevertheless, the global-equality characterization in Conjecture \texttt{00000008539} fails completely.

\begin{example}[Minimal Counterexample $B_2$]
\label{ex:b2}
Let $B_2$ be the Boolean lattice of subsets of $\{0, 1\}$, a geometric lattice of rank $n = 2$ with $4$ elements: $\emptyset$ (rank $0$), $\{0\}, \{1\}$ (rank $1$), and $\{0, 1\}$ (rank $2$). By Theorem~\ref{thm:bool-char}, $|\mu(\emptyset, S)| = 1$ for every $S \subseteq \{0, 1\}$, so
\[
  \bigl(w_0^+(B_2),\, w_1^+(B_2),\, w_2^+(B_2)\bigr) \;=\; \left(\binom{2}{0},\, \binom{2}{1},\, \binom{2}{2}\right) \;=\; (1, 2, 1).
\]
At $k = 1$, $(w_1^+(B_2))^2 = 2^2 = 4 \neq 1 = 1 \cdot 1 = w_0^+(B_2)\,w_2^+(B_2)$. More generally, on every Boolean lattice $B_n$ of rank $n \ge 2$, $w_k^+(B_n) = \binom{n}{k}$ satisfies strict inequality at every interior index $1 \le k \le n - 1$:
\[
  \frac{\binom{n}{k}^2}{\binom{n}{k-1}\binom{n}{k+1}} \;=\; \frac{k + 1}{k} \cdot \frac{n - k + 1}{n - k} \;>\; 1.
\]
\end{example}

Even more strikingly, we prove that \emph{no} geometric lattice of rank $\ge 2$ can attain global equality for $w_k^+(L)$.

\begin{theorem}[Complete Classification of Global Equality for $w_k^+(L)$]
\label{thm:whitney-class}
Let $L$ be any finite geometric lattice of rank $n = \mathrm{rank}(L)$.
\begin{enumerate}
  \item If $n \ge 2$, global log-concavity equality $(w_k^+(L))^2 = w_{k-1}^+(L)\,w_{k+1}^+(L)$ ($\forall 1 \le k \le n-1$) \textbf{never} holds.
  \item Consequently, global equality holds if and only if $n \le 1$ (where the interior index set $\{1,\dots,n-1\}$ is empty and $L \cong B_0$ or $B_1$).
  \item In particular, Conjecture \texttt{00000008539} under Convention I (asserting that global equality holds if and only if $L$ is Boolean) is \textbf{false}.
\end{enumerate}
\end{theorem}

\begin{proof}
Suppose $n \ge 2$ and $(w_k^+(L))^2 = w_{k-1}^+(L)\,w_{k+1}^+(L)$ for all $1 \le k \le n-1$. First, $w_0^+(L) = 1$. Second, let $a = w_1^+(L)$. Since $n \ge 2$, $\hatone \neq \hatzero$, so there is an atom $p_1 \le \hatone$; if $p_1$ were the only atom of $L$, atomisticity would give $\hatone \le p_1$ and $2 \le \rk(\hatone) \le \rk(p_1) = 1$, a contradiction. Thus $L$ has at least two distinct atoms $p_1 \neq p_2$, each with $|\mu(\hatzero, p_i)| \ge 1$, giving $a = w_1^+(L) \ge 2$.

By strong induction on $k \le n$, the recurrence $(w_k^+)^2 = w_{k-1}^+ w_{k+1}^+$ with $w_0^+ = 1$, $w_1^+ = a \ge 2$ yields $a^{k-1} w_{k+1}^+ = (a^k)^2 = a^{k-1} a^{k+1}$, and cancelling $a^{k-1} > 0$ gives $w_k^+(L) = a^k$ for all $0 \le k \le n$.

On the other hand, since $n \ge 2$, $\hatzero \neq \hatone$, so $(\mu * \zeta)(\hatzero, \hatone) = \delta(\hatzero, \hatone) = 0$. Partitioning by rank and applying Theorem~\ref{thm:rota} ($\mu(\hatzero, x) = (-1)^{\rk(x)} |\mu(\hatzero, x)|$) gives
\[
  0 \;=\; \sum_{x \in L} \mu(\hatzero, x) \;=\; \sum_{k=0}^n (-1)^k w_k^+(L) \;=\; \sum_{k=0}^n (-a)^k.
\]
Multiplying by $1 + a$ telescopes to $0 = (1 + a) \sum_{k=0}^n (-a)^k = 1 - (-a)^{n+1}$, forcing $a^{n+1} = |(-a)^{n+1}| = 1$, which is impossible since $a \ge 2$ and $n + 1 \ge 3$ imply $a^{n+1} \ge 8$. This proves~(1); (2) follows immediately since $\{k \in \N : 1 \le k \le n-1\} = \emptyset$ when $n \le 1$, and~(3) follows from Example~\ref{ex:b2}.
\end{proof}

\begin{remark}[Degenerate One-Way Implication]
By Theorem~\ref{thm:whitney-class}(2) and Lemma~\ref{lem:rk1-bool}, the bare material implication $\mathrm{GlobalEq}(w^+(L)) \Rightarrow \mathrm{IsBoolean}(L)$ holds formally (proved as theorem \texttt{whitney1Abs\_\allowbreak global\_\allowbreak eq\_\allowbreak imp\_\allowbreak boolean} in \texttt{lean/Main.lean}), \emph{solely} because $\mathrm{GlobalEq}(w^+(L))$ is false for all $n \ge 2$ and vacuous for $n \le 1$. However, in both English (\emph{``and equality holds globally only for Boolean lattices''}) and Chinese (\emph{``denghao jin zai buerge zhengti chengli''}), specifying the equality condition of an inequality with the modifier \emph{``globally''} presupposes that Boolean lattices of general rank actually attain equality across the non-empty interior range $1 \le k \le n-1$, which Example~\ref{ex:b2} refutes.
\end{remark}

\section{Conventions II and III: Elementwise, Averaged, and Normalized Sequences}

We now examine the literal elementwise reading of ``the absolute values of the M\"obius function'' (and its rank-averaged, rank-extremal, and binomial-normalized counterparts), where Boolean lattices $B_n$ \emph{do} attain global equality for every $n \ge 0$.

\begin{example}[Minimal Counterexample $M_3$ for Conventions IIA, IIB, IIC]
\label{ex:m3}
Let $M_3$ be the $5$-element diamond lattice $\{\hatzero, p_1, p_2, p_3, \hatone\}$ with three incomparable atoms $p_1, p_2, p_3$ satisfying $p_i \wedge p_j = \hatzero$ and $p_i \vee p_j = \hatone$ ($i \neq j$). Then $M_3 \cong U_{2,3} \cong \Pi_3$ is a rank-$2$ non-Boolean geometric lattice (distributivity fails: $p_1 \wedge (p_2 \vee p_3) = p_1 \neq \hatzero = (p_1 \wedge p_2) \vee (p_1 \wedge p_3)$). Its M\"obius values from $\hatzero$ are
\[
  \mu(\hatzero, \hatzero) = 1, \qquad \mu(\hatzero, p_1) = \mu(\hatzero, p_2) = \mu(\hatzero, p_3) = -1, \qquad \mu(\hatzero, \hatone) = -(1 - 3) = 2.
\]
Along \emph{every} maximal chain $\hatzero \lessdot p_i \lessdot \hatone$ of $M_3$, and likewise for the rank-average $\bar{\mu}_k(M_3) = w_k^+(M_3)/W_k(M_3)$ (giving $(1/1, 3/3, 2/1)$), rank-minimum, and rank-maximum, the resulting sequence is
\[
  (a_0, a_1, a_2) \;=\; (1, 1, 2),
\]
which \textbf{violates log-concavity} at $k = 1$:
\[
  a_1^2 \;=\; 1^2 \;=\; 1 \;<\; 2 \;=\; 1 \cdot 2 \;=\; a_0\,a_2.
\]
\end{example}

Theorem~\ref{thm:bool-char} upgrades Example~\ref{ex:m3} to a universal obstruction on \emph{every} non-Boolean geometric lattice.

\begin{theorem}[Universal Failure of Elementwise and Rank-Averaged Log-Concavity]
\label{thm:elem-universal}
Let $L$ be a finite geometric lattice of rank $n = \mathrm{rank}(L)$, and let $s : \{0,\dots,n\} \to L$ be any rank-selector ($\rk(s(k)) = k$ for $0 \le k \le n$, including every maximal chain $\hatzero = s(0) \lessdot s(1) \lessdot \dots \lessdot s(n) = \hatone$). Set $a_k = |\mu(\hatzero, s(k))|$ and $\bar{\mu}_k = w_k^+(L)/W_k(L)$.
\begin{enumerate}
  \item If $L$ is Boolean, then $a_k = 1$ and $\bar{\mu}_k = 1$ for all $0 \le k \le n$; hence both sequences have first term $1$, are log-concave, and satisfy global equality for every $n \ge 0$.
  \item If $L$ is \textbf{not} Boolean, then neither $(a_0,\dots,a_n)$ nor $(\bar{\mu}_0,\dots,\bar{\mu}_n)$ (nor the rank-minimum or rank-maximum sequence) is log-concave.
  \item Consequently, $(a_0,\dots,a_n)$ is log-concave if and only if $L$ is a Boolean lattice, and Conjecture \texttt{00000008539} is \textbf{false} under Conventions IIA, IIB, and IIC.
\end{enumerate}
\end{theorem}

\begin{proof}
Part~(1) is immediate from Theorem~\ref{thm:bool-char} ($|\mu(\hatzero, y)| = 1$ for all $y \in L$). For~(2), suppose $L$ is not Boolean. By Lemma~\ref{lem:rk1-bool}, $n \ge 2$, and by Theorems~\ref{thm:rota} and~\ref{thm:bool-char}, $|\mu(\hatzero, \hatone)| \ge 2$. Since $\rk(s(0)) = 0$, $s(0) = \hatzero$ and $a_0 = |\mu(\hatzero, \hatzero)| = 1$. Since $\rk(s(1)) = 1$, $s(1)$ is an atom ($\hatzero \lessdot s(1)$), so $\mu(\hatzero, s(1)) = -1$ by Proposition~\ref{prop:mobius} and $a_1 = 1$; similarly, every rank-$1$ element has $|\mu| = 1$, so $\bar{\mu}_0 = \bar{\mu}_1 = 1$. Finally, $\rk(s(n)) = n = \rk(\hatone)$ forces $s(n) = \hatone$, so $a_n = \bar{\mu}_n = |\mu(\hatzero, \hatone)| \ge 2$.

However, if any sequence $(c_0, c_1, \dots, c_n)$ of real numbers $\ge 1$ satisfies $c_0 = 1$, $c_1 = 1$, and $c_k^2 \ge c_{k-1} c_{k+1}$ for $1 \le k \le n - 1$, strong induction on $k$ yields $c_k = 1$ for all $0 \le k \le n$: assuming $c_{k-1} = c_k = 1$, the inequality $1 \cdot c_{k+1} = c_{k-1} c_{k+1} \le c_k^2 = 1$ combined with $c_{k+1} \ge 1$ forces $c_{k+1} = 1$. Applying this to $a_k$ and $\bar{\mu}_k$ (and rank-min/max, which also satisfy $c_0 = c_1 = 1$ and $c_n = |\mu(\hatzero,\hatone)| \ge 2$) contradicts $c_n \ge 2$.
\end{proof}

\begin{theorem}[Disproof of Convention III --- Binomial-Normalized Whitney Numbers]
\label{thm:ulc}
Let $U_{3,4}$ be the rank-$3$ uniform geometric lattice on $\{0,1,2,3\}$, whose $12$ flats consist of all subsets of cardinality $0, 1, 2$ together with $\{0,1,2,3\}$ at rank $3$. Then $U_{3,4}$ is a finite geometric lattice with unsigned Whitney sequence
\[
  \bigl(w_0^+(U_{3,4}),\, w_1^+(U_{3,4}),\, w_2^+(U_{3,4}),\, w_3^+(U_{3,4})\bigr) \;=\; (1, 4, 6, 3).
\]
Its binomial-normalized sequence $b_k(U_{3,4}) = w_k^+(U_{3,4}) / \binom{3}{k} = \bigl(1, \frac{4}{3}, 2, 3\bigr)$ \textbf{violates log-concavity} at $k = 1$:
\[
  b_1(U_{3,4})^2 \;=\; \left(\frac{4}{3}\right)^2 \;=\; \frac{16}{9} \;<\; 2 \;=\; 1 \cdot 2 \;=\; b_0(U_{3,4})\,b_2(U_{3,4}).
\]
\end{theorem}

\begin{proof}
For $k \in \{0,1,2\}$, every $k$-element subset $F \subseteq \{0,1,2,3\}$ has interval $[\emptyset, F] \cong B_k$, so $|\mu(\emptyset, F)| = 1$ and $w_k^+(U_{3,4}) = \binom{4}{k}$, giving $w_0^+ = 1$, $w_1^+ = 4$, $w_2^+ = 6$. Since $\sum_{x \in U_{3,4}} \mu(\emptyset, x) = 0$, we obtain $\mu(\emptyset, \{0,1,2,3\}) = -(1 - 4 + 6) = -3$, so $w_3^+(U_{3,4}) = 3$. Dividing by $(\binom{3}{0}, \binom{3}{1}, \binom{3}{2}, \binom{3}{3}) = (1, 3, 3, 1)$ gives $(1, 4/3, 2, 3)$, where $(4/3)^2 = 16/9 < 2$.
\end{proof}

\section{Formal Verification in Core Lean 4 and Independent Python Audit}

All results in Sections~2--5 are formally proved in \texttt{lean/Main.lean} using self-contained Core Lean~4 (toolchain \texttt{leanprover/lean4:v4.19.0}, zero external dependencies, compiled with \texttt{-DwarningAsError=true}, zero \texttt{sorry}, zero \texttt{native\_decide}, and zero custom axioms). Table~\ref{tab:lean} maps the paper's definitions and theorems to their Lean declarations.

In addition, \texttt{verify.py} independently constructs $14$ finite geometric lattices ($B_0$--$B_4$, $M_3$--$M_5$, $U_{3,4}$, $U_{3,5}$, $U_{4,5}$, the partition lattice $\Pi_4$, the Fano plane $\mathrm{PG}(2,2)$, and $M_3 \times B_1$) from flat closure systems and verifies all incidence-algebra, lattice, sign, log-concavity, and counterexample assertions.

\newpage
\begin{table}[!t]
\centering
\small
\renewcommand{\arraystretch}{1.24}
\begin{tabular}{@{}ll@{}}
\hline
\textbf{Mathematical Result} & \textbf{Lean 4 Declaration in \texttt{lean/Main.lean}} \\
\hline
Incidence algebra two-sided inverse $\mu * \zeta = \delta = \zeta * \mu$ & \texttt{IncidenceAlg.mobius\_is\_mobius} \\
Uniqueness of $\mu$ in $I(P, \Z)$ & \texttt{IncidenceAlg.mobius\_unique} \\
Uniqueness of geometric lattice rank $\rk$ & \texttt{GeometricLattice.rk\_unique} \\
Weisner's identity (Theorem~\ref{thm:weisner}) & \texttt{weisner\_thm} \\
Rota's strict sign theorem (Theorem~\ref{thm:rota}) & \texttt{rota\_sign\_thm} \\
Boolean $\Leftrightarrow |\mu(\hatzero,\hatone)|=1$ (Theorem~\ref{thm:bool-char}) & \texttt{isBoolean\_iff\_natAbs\_mu\_top\_eq\_one} \\
No global equality for $w_k^+$ when $\rk \ge 2$ (Theorem~\ref{thm:whitney-class}(1)) & \texttt{whitney1Abs\_no\_global\_eq\_of\_rank\_ge\_two} \\
Global equality for $w_k^+ \Leftrightarrow \rk \le 1$ (Theorem~\ref{thm:whitney-class}(2)) & \texttt{whitney1Abs\_global\_eq\_iff\_rank\_le\_one} \\
Disproof under Convention I on $B_2$ (Theorem~\ref{thm:whitney-class}(3)) & \texttt{disproof\_convention\_I\_whitney\_iff} \\
Universal chain log-concavity failure (Theorem~\ref{thm:elem-universal}) & \texttt{nonBoolean\_rankSelector\_not\_logConcave} \\
Chain log-concavity $\Leftrightarrow$ Boolean (Theorem~\ref{thm:elem-universal}(3)) & \texttt{rankSelector\_logConcave\_iff\_boolean} \\
Universal rank-average failure (Theorem~\ref{thm:elem-universal}(2)) & \texttt{nonBoolean\_avg\_not\_logConcave} \\
Disproofs under Conventions IIA \& IIB on $M_3$ & \texttt{disproof\_convention\_IIA\_chain}, \texttt{\_IIB\_avg} \\
Disproof under Convention III on $U_{3,4}$ (Theorem~\ref{thm:ulc}) & \texttt{disproof\_convention\_III\_ultra} \\
Non-vacuity witnesses ($B_2$, $M_3$, $U_{3,4}$) & \texttt{positive\_non\_vacuity\_witnesses} \\
\hline
\end{tabular}
\caption{Correspondence between mathematical theorems and Core Lean 4 declarations.}
\label{tab:lean}
\end{table}

\begin{thebibliography}{9}

\bibitem{AHK2018}
K.~Adiprasito, J.~Huh, and E.~Katz,
\emph{Hodge theory for combinatorial geometries},
Ann.\ of Math.\ (2) \textbf{188} (2018), no.~2, 381--452.

\bibitem{Rota1964}
G.-C.~Rota,
\emph{On the foundations of combinatorial theory I. Theory of M\"obius functions},
Z.\ Wahrscheinlichkeitstheorie und Verw.\ Gebiete \textbf{2} (1964), 340--368.

\bibitem{StanleyEC1}
R.~P.~Stanley,
\emph{Enumerative Combinatorics, Volume 1}, 2nd ed.,
Cambridge Studies in Advanced Mathematics, vol.~49, Cambridge Univ.\ Press, 2012.

\end{thebibliography}

\end{document}
